% ==========================================================================
%  Chapter 104 -- Architecture of the Branch-Selection Extension (diagrams)
% ==========================================================================

\chapter{Architecture of the Branch-Selection Extension}
\label{ch:extension-architecture}

\paragraph{Normative status.}
Design record for the extension of
Caps.~\ref{ch:branch-selection-reversals}--\ref{ch:ca-tuner-solver}. It
fixes the type architecture (concepts, policies, mixins) and the control
flow (the detect--restore--retry loop) so the driver integration of
Phase~I is a wiring exercise, not a design one. No new runtime behaviour
is introduced here.

\section{Design rule: concepts first, CRTP only as a static mixin}
\label{sec:ext-design-rule}

The extension obeys one architectural rule, applied without exception:

\begin{quote}
\emph{Contracts are expressed as C++20 \texttt{concept}s. CRTP is used
\textbf{only} as a ``static mixin'': to \textbf{inject} operations into a
strong type, never to express an interface or to dispatch statically.}
\end{quote}

The distinction is between two jobs that inheritance is often
(mis)used for:
%
\begin{description}
  \item[Contract / static duck-typing.] ``This template argument must
        provide \texttt{residual}, \texttt{tangent}, \texttt{accept}.''
        This is a \texttt{concept}. It constrains without a base class, it
        composes (a type can model several concepts), it gives readable
        errors, and it never forces an inheritance edge. Every seam in the
        extension --- \texttt{ContinuationBackend},
        \texttt{RegularizationPolicy}, \texttt{DeflationPolicy},
        \texttt{EnergyContinuationBackend}, \texttt{SearchSpace},
        \texttt{ObjectiveFunction}, \texttt{KnowledgeSource} --- is a
        concept, and each model carries a \texttt{static\_assert} that it
        satisfies its concept.
  \item[Operation injection into a strong type.] ``This strong type
        carries data and identity; graft a family of operations onto it
        that are written once and reused across many such types.'' Here a
        \emph{static mixin} via CRTP is the right tool: the mixin is a
        \texttt{template<class Derived> struct Mixin} whose methods
        \texttt{static\_cast} to \texttt{Derived} to reach the type's
        data, and the strong type derives \texttt{: Mixin<Strong>}. The
        operations live in the mixin (reusable, extensible), the state
        lives in the strong type. The extension does not \emph{need} a
        mixin yet --- its policies are small structs of traits --- but the
        seam is reserved: e.g. a future strong \texttt{Genome} type would
        receive its arithmetic (clamp, blend, distance) from a
        \texttt{GenomeOps<Genome>} mixin rather than re-implementing them.
\end{description}

The rule is why the whole extension has \emph{zero} CRTP today: nothing
yet injects operations into a strong type; everything is a contract, and
contracts are concepts.

\section{Type architecture --- solver side (P1, P2)}

Figure~\ref{fig:ext-solver-types} is the concept/class graph of the two
continuation solvers. Dashed boxes are concepts; solid boxes are class
templates or their policy arguments. A dashed edge reads ``models''; a
solid edge reads ``is a template parameter of''.

\begin{figure}[htbp]
\centering
\begin{tikzpicture}[
    concept/.style={draw, dashed, rounded corners=2pt, fill=blue!6,
                    font=\scriptsize\ttfamily, align=center, inner sep=3.5pt},
    klass/.style={draw, rounded corners=2pt, fill=green!8,
                  font=\scriptsize\ttfamily, align=center, inner sep=3.5pt},
    policy/.style={draw, rounded corners=2pt, fill=orange!10,
                   font=\scriptsize\ttfamily, align=center, inner sep=3.5pt},
    models/.style={-{Stealth[length=4pt]}, dashed, gray},
    tmpl/.style={-{Stealth[length=5pt]}, thick}
  ]
  % concept row
  \node[concept] (cb)  at (0,5)    {ContinuationBackend};
  \node[concept] (rp)  at (4.6,5)  {RegularizationPolicy};
  \node[concept] (dp)  at (9.2,5)  {DeflationPolicy};
  \node[concept] (ecb) at (0,3.4)  {EnergyContinuationBackend};
  % classes
  \node[klass] (rnc) at (5.0,3.0)
     {RegularizedNewtonContinuation\\<Backend, Reg, Defl>};
  \node[klass] (tao) at (0.2,1.6) {TaoEnergyContinuation<Backend>};
  % policies
  \node[policy] (lm)  at (3.4,0.4) {LevenbergMarquardt};
  \node[policy] (pn)  at (6.0,0.4) {PureNewton};
  \node[policy] (nd)  at (8.4,1.3) {NoDeflation};
  \node[policy] (smd) at (9.6,0.2) {ShermanMorrison\\Deflation};

  % models edges
  \draw[models] (ecb) -- (cb) node[midway,left,font=\tiny\rmfamily]{refines};
  \draw[models] (lm)  -- (rp);
  \draw[models] (pn)  -- (rp);
  \draw[models] (nd)  -- (dp);
  \draw[models] (smd) -- (dp);
  % template-parameter edges into the solver
  \draw[tmpl] (cb)  -- (rnc);
  \draw[tmpl] (rp)  -- (rnc);
  \draw[tmpl] (dp)  -- (rnc);
  \draw[tmpl] (ecb) -- (tao);
\end{tikzpicture}
\caption[Concept/class graph of the continuation solvers.]{Concept/class
graph of the branch-selection solvers. Every seam is a concept (dashed);
policies are swapped as template arguments. \texttt{NoDeflation} is the
default third argument, so the un-deflated solver is one instantiation of
the same template. \texttt{EnergyContinuationBackend} \emph{refines}
\texttt{ContinuationBackend} by adding \texttt{incremental\_energy}.}
\label{fig:ext-solver-types}
\end{figure}

The reading of Figure~\ref{fig:ext-solver-types}: the solver
\texttt{RegularizedNewtonContinuation<Backend,Reg,Defl>} is parameterised
by three \emph{independent} concepts. Swapping \texttt{Defl} from
\texttt{NoDeflation} to \texttt{ShermanMorrisonDeflation} turns deflation
on with no other change; because the default is \texttt{NoDeflation} and
all deflation code sits behind \texttt{if constexpr (Defl::enabled)}, the
default instantiation is byte-identical to the pre-extension solver
(Cap.~\ref{ch:deflated-continuation}). \texttt{TaoEnergyContinuation}
reuses the same backend contract, refined with one method.

\section{Type architecture --- metaheuristics side (M, P3)}

Figure~\ref{fig:ext-algo-types} is the graph of the optimisation module.
The same rule holds: \texttt{SearchSpace}, \texttt{ObjectiveFunction} and
\texttt{KnowledgeSource} are concepts; \texttt{CulturalAlgorithm} and
\texttt{BeliefSpace} are class templates parameterised by them.

\begin{figure}[htbp]
\centering
\begin{tikzpicture}[
    concept/.style={draw, dashed, rounded corners=2pt, fill=blue!6,
                    font=\scriptsize\ttfamily, align=center, inner sep=3.5pt},
    klass/.style={draw, rounded corners=2pt, fill=green!8,
                  font=\scriptsize\ttfamily, align=center, inner sep=3.5pt},
    policy/.style={draw, rounded corners=2pt, fill=orange!10,
                   font=\scriptsize\ttfamily, align=center, inner sep=3.5pt},
    models/.style={-{Stealth[length=4pt]}, dashed, gray},
    tmpl/.style={-{Stealth[length=5pt]}, thick},
    comp/.style={-{Diamond[open,length=6pt]}, thick}
  ]
  \node[concept] (ss)  at (0,4.6)   {SearchSpace};
  \node[concept] (of)  at (4.4,4.6) {ObjectiveFunction};
  \node[concept] (ks)  at (9.0,4.6) {KnowledgeSource};
  \node[klass]   (ca)  at (2.4,2.7) {CulturalAlgorithm\\<Space, KS...>};
  \node[klass]   (bs)  at (7.4,2.7) {BeliefSpace<KS...>};
  \node[policy]  (bss) at (0,1.2)   {BoundedSearchSpace};
  \node[policy]  (nk)  at (7.0,1.1) {NormativeKnowledge};
  \node[policy]  (sk)  at (10.2,1.1){SituationalKnowledge};

  \draw[models] (bss) -- (ss);
  \draw[models] (nk)  -- (ks);
  \draw[models] (sk)  -- (ks);
  \draw[tmpl] (ss) -- (ca);
  \draw[tmpl] (ks) -- (ca);
  \draw[tmpl] (ks) -- (bs);
  \draw[comp] (ca) -- (bs) node[midway,above,font=\tiny\rmfamily]{owns};
  \draw[tmpl] (of.south) .. controls (4.4,3.4) .. (ca.north east)
        node[near start,right,font=\tiny\rmfamily]{maximize(\,)};
\end{tikzpicture}
\caption[Concept/class graph of the metaheuristics module.]{Concept/class
graph of \texttt{src/algorithms/}. \texttt{CulturalAlgorithm} is
constrained by \texttt{SearchSpace} and a pack of \texttt{KnowledgeSource}s
and owns a \texttt{BeliefSpace} over the same pack; the objective enters
through the \texttt{maximize} method (an \texttt{ObjectiveFunction}).
Adding a knowledge source is a template-argument change, not an
inheritance edit --- the belief space fans over the pack with a fold
expression.}
\label{fig:ext-algo-types}
\end{figure}

\section{Control flow --- the detect--restore--retry loop (Phase I)}

The reversal remedy is reactive: a spurious step is only recognisable
\emph{after} it commits, because the crack ratchet has advanced by then.
The driver therefore brackets each reversal step with a checkpoint and
retries on detection, using the Phase-P1 API
(\texttt{capture\_state}, \texttt{register\_spurious\_root},
\texttt{restore\_state}). Figure~\ref{fig:ext-retry-flow} is the process.

\begin{figure}[htbp]
\centering
\begin{tikzpicture}[
    node distance=6mm,
    pnode/.style={draw, rounded corners=2pt, fill=green!7, font=\scriptsize,
                 align=center, inner sep=4pt, minimum width=4.4cm},
    dec/.style={draw, rounded corners=8pt, fill=orange!12, font=\scriptsize,
                align=center, inner sep=4pt, minimum width=2.4cm},
    arr/.style={-{Stealth[length=5pt]}, thick},
    lbl/.style={font=\tiny, fill=white, inner sep=1pt}
  ]
  \node[pnode] (chk)  {checkpoint: \texttt{capture\_state()}};
  \node[pnode, below=of chk] (adv) {\texttt{advance\_to(p)} (LM step)};
  \node[dec, below=8mm of adv] (conv) {converged\\\&\ no spike?};
  \node[pnode, right=22mm of conv, fill=blue!7] (commit) {commit step,\\advance to next \(p\)};
  \node[dec, below=10mm of conv] (retry) {retries\\left?};
  \node[pnode, below=9mm of retry] (reg)
       {\texttt{register\_spurious\_root(u)}\\\texttt{restore\_state()}, tighten \(\mu\)};
  \node[pnode, left=20mm of retry, fill=red!7] (accept)
       {accept regularized\\(as baseline)};

  \draw[arr] (chk) -- (adv);
  \draw[arr] (adv) -- (conv);
  \draw[arr] (conv) -- node[lbl]{yes} (commit);
  \draw[arr] (conv) -- node[lbl]{no} (retry);
  \draw[arr] (retry) -- node[lbl]{no} (accept);
  \draw[arr] (retry) -- node[lbl]{yes} (reg);
  \draw[arr] (reg.west) -- ++(-1.2,0) |- (adv.west)
        node[lbl, pos=0.25]{re-solve, deflated};
\end{tikzpicture}
\caption[Detect--restore--retry loop at a reversal.]{The Phase-I
detect--restore--retry loop. A reversal step is checkpointed, advanced,
and audited; a converged spurious step (high force spike) is rolled back,
its state registered as a deflation root, and re-solved so the deflated
continuation escapes the spurious basin. On exhausting the retry budget
the step is accepted \emph{as the baseline would have}, so the gate-off
behaviour is preserved. The rounded amber nodes are decisions.}
\label{fig:ext-retry-flow}
\end{figure}

Two invariants make the loop safe. First, the retry budget defaults to
zero, so with the gate off the loop degenerates to a single
\texttt{advance\_to} and the baseline path (Cap.~\ref{ch:branch-selection-reversals}).
Second, the roots registered during a step are cleared when the step is
committed, because the multi-equilibrium is local to a fixed control value
--- a root that repelled the iteration at one reversal must not repel the
physical branch at the next.

\section{Where the CA-tuner attaches}

The Phase-P3 meta-tuner (Cap.~\ref{ch:ca-tuner-solver}) wraps the same
checkpoint: it captures the state before the first reversal, and its
objective decodes a genome into a config, restores, replays the window,
scores, and restores again. The tuner is thus a \emph{second} consumer of
the Phase-P1 checkpoint API, orthogonal to the retry loop --- the retry
loop selects a branch \emph{within} a fixed config; the tuner selects a
\emph{config}. Both leave the default build byte-identical because both
are entered only under their environment gates.
